\chapter{Peripheral Artery Bypass}

\section*{Introduction \& Clinical Indications}
Peripheral artery bypass is a significant surgical intervention aimed at restoring blood flow to the lower extremities in patients afflicted with severe peripheral artery disease (PAD). This condition arises from the narrowing or occlusion of arteries, predominantly in the legs, resulting in diminished blood supply and subsequent tissue ischemia. The bypass procedure entails the creation of an alternative pathway for blood to circumvent an obstructed or stenotic segment of an artery, utilizing either an autologous vein graft, typically harvested from the great saphenous vein, or a synthetic conduit. Common variations of this procedure include aortofemoral, femoropopliteal, and femorotibial bypasses, named according to the arteries involved. The clinical significance of this revascularization lies in its ability to alleviate debilitating symptoms such as claudication, ischemic rest pain, and non-healing ulcers or gangrene, ultimately aiming to prevent limb loss, enhance functional capacity, and improve the overall quality of life for affected patients.

\begin{itemize}
  \item Lower-extremity revascularization surgery
  \item Leg bypass surgery
  \item Femoropopliteal bypass
  \item Aortofemoral bypass
  \item Femorotibial bypass
  \item Distal bypass surgery
  \item Iliofemoral bypass
\end{itemize}

The fundamental principle underlying peripheral artery bypass surgery is to bypass an arterial obstruction by establishing a new conduit that connects a healthy artery proximal to the blockage (inflow) to a healthy artery distal to the blockage (outflow). This re-establishment of pulsatile blood flow to the ischemic limb enhances oxygen and nutrient delivery to the tissues, promotes wound healing, and alleviates ischemic symptoms. The choice of conduit is critical; autologous saphenous vein is generally preferred due to its superior long-term patency, particularly in distal bypasses, owing to its biological compatibility, resistance to infection, and ability to remodel. In cases where suitable vein is unavailable, synthetic grafts made of materials such as polytetrafluoroethylene (PTFE) or Dacron are utilized, especially for larger diameter vessels and less distal bypasses, although these carry a higher risk of infection and lower patency rates in smaller, more distal vessels.

The surgical rationale focuses on identifying a suitable inflow vessel with adequate pressure and flow, as well as a patent outflow vessel or segment of artery distal to the disease that can accept the graft and effectively perfuse the limb. Technical goals include achieving tension-free anastomoses to prevent kinking or tearing, ensuring proper graft length and diameter to minimize turbulence and optimize flow dynamics, and maintaining strict hemostasis throughout the procedure to prevent hematoma formation. Ultimately, the aim is to restore sufficient perfusion to heal tissue loss, relieve ischemic pain, improve functional capacity, and prevent major amputation, thereby enhancing the patient's quality of life.

The evolution of peripheral artery bypass surgery marks a pivotal advancement in vascular medicine, fundamentally transforming the management of limb-threatening ischemia. Early attempts at arterial reconstruction in the late 19th and early 20th centuries were largely unsuccessful due to complications such as thrombosis, infection, and the absence of suitable graft materials. A landmark moment occurred in 1948 when Jean Kunlin, a French surgeon, successfully performed the first femoropopliteal bypass using an autologous saphenous vein graft. This groundbreaking technique demonstrated the feasibility of utilizing a patient's own vein as a bypass conduit, laying the groundwork for modern revascularization.

Following Kunlin's innovation, the field rapidly progressed. In the 1950s, pioneering surgeons such as Michael DeBakey and Denton Cooley refined the use of synthetic materials, including Dacron and later PTFE, for arterial reconstruction, particularly for larger vessels like the aorta and iliac arteries in aortofemoral bypasses. The development of finer sutures, improved surgical instruments, and a deeper understanding of hemodynamics and coagulation further enhanced surgical techniques. Over the subsequent decades, the focus shifted towards improving long-term patency, reducing complications, and extending bypass options to more distal vessels, including tibial and pedal arteries, to salvage limbs with critical ischemia. The introduction of duplex ultrasound for graft surveillance and advanced imaging techniques for preoperative planning have further optimized patient selection and surgical outcomes, making bypass a durable option for complex peripheral artery disease.

Peripheral artery bypass is indicated for patients with severe peripheral artery disease (PAD) when less invasive treatments have failed or are deemed unsuitable. Specific clinical indications include:

\begin{itemize}
  \item Critical limb ischemia (CLI) characterized by persistent ischemic rest pain requiring regular analgesia for more than two weeks, often classified as Rutherford Category 4.
  \item Non-healing arterial ulcers or gangrene of the foot or toes due to insufficient blood flow, corresponding to Rutherford Categories 5 and 6.
  \item Lifestyle-limiting claudication that significantly impairs daily activities and quality of life, despite optimal medical therapy and supervised exercise programs, typically Rutherford Category 3.
  \item Acute limb ischemia (ALI) when endovascular options are not feasible or have failed, and there is an immediate threat of limb loss.
  \item Failed previous endovascular interventions (e.g., angioplasty, stenting) that have resulted in re-occlusion or inadequate revascularization.
  \item Anatomical considerations, such as long-segment occlusions, heavily calcified arteries, or complex lesions that are not amenable to endovascular techniques.
  \item Progressive ischemic neuropathy causing severe pain or motor dysfunction, unresponsive to conservative management.
  \item Symptomatic popliteal artery aneurysm, particularly if it is thrombosed or embolizing.
\end{itemize}

Peripheral artery bypass can function as either a primary or secondary revascularization procedure, contingent upon the patient's clinical presentation, the anatomical extent of disease, and prior interventions. It is considered a primary procedure when the patient presents with critical limb ischemia (CLI) and the arterial anatomy is complex, involving long-segment occlusions, extensive calcification, or multiple levels of disease that are not amenable to endovascular approaches. In such cases, bypass surgery often provides a more durable and comprehensive revascularization, particularly for infrainguinal disease. Conversely, bypass surgery frequently serves as a secondary or salvage procedure after initial endovascular attempts (such as angioplasty or stenting) have failed to achieve or maintain adequate perfusion, or when re-stenosis occurs. It may also be a secondary option for patients whose claudication persists despite maximal medical therapy and supervised exercise, and who have failed endovascular treatment. The decision to proceed with bypass as primary or secondary therapy is made after a comprehensive evaluation by a multidisciplinary team, weighing the risks and benefits of all available revascularization strategies, considering factors such as limb threat, patient comorbidities, and expected long-term patency.

\section*{Relevant Surgical Anatomy}
A comprehensive understanding of the relevant surgical anatomy is crucial for the successful execution of peripheral artery bypass surgery. The arterial system involved typically includes the aorta, common iliac arteries, external iliac arteries, common femoral arteries, superficial femoral arteries, profunda femoris arteries, popliteal arteries, and the trifurcation vessels, which consist of the anterior tibial, posterior tibial, and peroneal arteries. The inflow vessel must be free of significant disease to ensure adequate pressure and flow, while the outflow vessel distal to the occlusion must be patent to accept the graft and adequately perfuse the limb.

Key anatomical landmarks include the inguinal ligament, which separates the external iliac artery from the common femoral artery, and the adductor canal (Hunter's canal), where the superficial femoral artery transitions into the popliteal artery. For autologous vein grafts, the great saphenous vein (GSV) is the conduit of choice due to its excellent long-term patency. The GSV originates from the dorsal venous arch of the foot, ascends anterior to the medial malleolus, and runs along the medial aspect of the calf and thigh, draining into the common femoral vein at the saphenofemoral junction. Important nerve relations include the saphenous nerve, which accompanies the GSV, and the common peroneal nerve, which is at risk during dissection around the popliteal fossa. Lymphatic drainage of the lower limb primarily follows the venous system to the inguinal lymph nodes, which can be disrupted during groin dissection, leading to complications such as lymphocele or lymphedema.

\section*{Preoperative Workup \& Preparation}
Thorough preoperative preparation is essential to optimize patient outcomes and minimize risks associated with peripheral artery bypass surgery. This typically involves a comprehensive medical evaluation and specific interventions to assess and mitigate cardiovascular risk, given the systemic nature of atherosclerosis.

Key aspects of preoperative workup and preparation include:

\begin{itemize}
  \item \textbf{Diagnostic Imaging:} Preoperative arterial mapping is crucial, usually performed with computed tomography angiography (CTA), magnetic resonance angiography (MRA), or conventional angiography, to precisely delineate the extent of arterial disease, identify suitable inflow and outflow vessels, and plan the graft trajectory and length. Duplex ultrasound may also be used to map the saphenous vein for suitability as a conduit.
  
  \item \textbf{Cardiac and Pulmonary Clearance:} Given the high prevalence of concomitant coronary artery disease and chronic obstructive pulmonary disease in PAD patients, a thorough cardiac workup (ECG, echocardiogram, stress testing, cardiology consultation) and pulmonary assessment (pulmonary function tests, smoking cessation counseling) are often required to ensure the patient can tolerate the physiological stress of surgery.
  
  \item \textbf{Laboratory Tests:} Routine blood tests include a complete blood count (CBC), electrolyte panel, renal and liver function tests, coagulation profile (PT/INR, PTT), and blood typing and cross-matching for potential transfusion. Glycosylated hemoglobin (HbA1c) is important for diabetic patients.
  
  \item \textbf{Medication Management:} Antiplatelet medications (e.g., aspirin, clopidogrel) are often continued, but anticoagulants (e.g., warfarin) may need to be bridged or temporarily discontinued per protocol. Beta-blockers and statins are typically continued. Diabetic patients require strict glucose control.
  
  \item \textbf{NPO Status:} Patients are instructed to be NPO (nil per os) for a specified period (typically 6-8 hours) before surgery to reduce the risk of aspiration during anesthesia.
  
  \item \textbf{Antibiotic Prophylaxis:} Intravenous broad-spectrum antibiotics (e.g., cefazolin) are administered shortly before the incision to prevent surgical site infection.
  
  \item \textbf{Informed Consent:} A detailed discussion with the surgeon regarding the procedure, potential benefits, risks, and alternatives is conducted, and the patient provides informed consent.
  
  \item \textbf{Nutritional Optimization:} Addressing any nutritional deficiencies can aid in wound healing and overall recovery, especially in patients with critical limb ischemia and chronic wounds.
\end{itemize}

Peripheral artery bypass surgery, while life-saving for many, carries specific contraindications and requires careful precautions to ensure patient safety and optimize outcomes.

\textbf{Absolute Contraindications:}
\begin{itemize}
  \item \textbf{Unreconstructable Distal Runoff:} The absence of a suitable, patent artery distal to the occlusion to connect the graft to, which would render the bypass futile due to immediate graft thrombosis.
  
  \item \textbf{Severe Systemic Illness:} Patients with end-stage heart failure (NYHA Class IV), severe uncontrolled sepsis, or other critical comorbidities that make the risks of major surgery prohibitive and outweigh any potential benefits of limb salvage.
  
  \item \textbf{Active Infection at Incision Site:} An ongoing infection in the planned surgical field, particularly in the groin, significantly increases the risk of graft infection, a devastating complication that often necessitates graft removal and limb amputation.
  
  \item \textbf{Inadequate Inflow Artery:} An inflow artery that is too diseased, stenotic, or occluded to provide sufficient blood pressure and flow to the bypass graft, leading to early graft failure.
  
  \item \textbf{Non-Ambulatory Status with Limited Life Expectancy:} For patients who are bed-bound with a very short life expectancy (e.g., less than 6 months), the morbidity and mortality of major surgery may not be justified, especially if limb salvage does not improve functional status or quality of life.
\end{itemize}

\textbf{Relative Contraindications and Precautions:}
\begin{itemize}
  \item \textbf{Poor Quality Autologous Vein Conduit:} If the saphenous vein is diseased, small in diameter (<3 mm), or previously harvested, synthetic grafts may be necessary, which have lower long-term patency rates, especially in distal bypasses. Alternative veins (e.g., arm veins) may be considered.
  
  \item \textbf{Heavy Arterial Calcification:} Extensive calcification of the native arteries can make clamping, suturing, and achieving a patent anastomosis technically challenging, increasing the risk of intimal dissection or anastomotic leak.
  
  \item \textbf{Severe Obesity (BMI >40 kg/m\textasciicircum2):} Increases surgical complexity, prolongs operative time, and significantly elevates wound complication rates (infection, dehiscence, lymphocele) and anesthetic risks.
  
  \item \textbf{Active Smoking:} Significantly impairs wound healing, accelerates atherosclerosis, and compromises long-term graft patency. Patients are strongly advised to cease smoking preoperatively, ideally for several weeks.
  
  \item \textbf{Uncontrolled Diabetes Mellitus (HbA1c >8\%):} Increases infection risk, impairs wound healing, and accelerates atherosclerosis, potentially compromising long-term graft patency. Strict glucose control is paramount.
  
  \item \textbf{End-Stage Renal Disease (ESRD):} Associated with increased cardiovascular risk, poorer graft outcomes, and higher rates of wound complications and infection.
  
  \item \textbf{Patient Non-compliance:} A patient unwilling or unable to adhere to postoperative medication regimens (e.g., antiplatelets, statins) or lifestyle modifications may have poorer long-term outcomes and higher rates of graft failure.
\end{itemize}

\section*{Surgical Technique \& Procedure}
Peripheral artery bypass surgery is a complex procedure performed under strict sterile conditions, typically requiring general anesthesia, although regional anesthesia (spinal or epidural) may be utilized in select cases. The patient is positioned supine, often with the operative leg slightly abducted and externally rotated (frog-leg position) to facilitate access to the groin, thigh, and calf.

The procedure generally follows these detailed steps:

\begin{enumerate}
  \item \textbf{Incision and Exposure:} Incisions are strategically made to expose the inflow artery (e.g., common femoral artery, aorta, or iliac artery) and the outflow artery (e.g., superficial femoral, popliteal, or tibial artery). For autologous vein grafts, a separate incision along the medial thigh and calf or an endoscopic technique is used to harvest the great saphenous vein from the same or contralateral leg, ensuring all venous branches are ligated.
  
  \item \textbf{Systemic Anticoagulation:} Once the arteries are exposed and prepared for clamping, the patient is systemically heparinized (typically 50-100 units/kg) to achieve an activated clotting time (ACT) of 250-300 seconds, preventing thrombosis during arterial clamping.
  
  \item \textbf{Arterial Clamping:} Vascular clamps are carefully applied to the inflow and outflow arteries to temporarily stop blood flow to the segment where the graft will be connected. This creates a bloodless field for precise anastomosis.
  
  \item \textbf{Proximal Anastomosis Creation:} The proximal anastomosis is created first, typically an end-to-side connection between the graft and the inflow artery, using fine non-absorbable monofilament sutures (e.g., 6-0 or 7-0 polypropylene). The arteriotomy is carefully sized to match the graft diameter.
  
  \item \textbf{Graft Tunneling:} A subcutaneous or subfascial tunnel is meticulously created using a tunneling device to route the graft from the proximal anastomosis to the distal anastomosis site. Care is taken to ensure the graft is not kinked, twisted, or under excessive tension, and that it avoids compression by surrounding structures.
  
  \item \textbf{Distal Anastomosis Creation:} The distal anastomosis is then created, usually end-to-side, connecting the graft to the outflow artery using similar fine non-absorbable sutures. This connection is critical for establishing flow to the distal limb.
  
  \item \textbf{Flow Restoration and Assessment:} After both anastomoses are complete, the clamps are carefully removed, typically in a sequence (distal outflow, then proximal inflow) to flush any debris. Blood flow is restored through the graft. Intraoperative assessment of graft patency and flow is performed using Doppler ultrasound, pulse palpation, or completion angiography to confirm technical success and identify any issues.
  
  \item \textbf{Hemostasis and Closure:} Meticulous hemostasis is achieved by ligating or cauterizing any bleeding vessels. Surgical drains (e.g., closed suction drains) may be placed in the groin or thigh to prevent hematoma or seroma formation. The incisions are then closed in layers, typically with absorbable sutures for deeper layers and non-absorbable sutures or staples for the skin.
\end{enumerate}

\section*{Complications \& Risks}
Peripheral artery bypass is a major surgical procedure with potential risks, which can be broadly categorized as general surgical risks and procedure-specific complications.

\textbf{General Surgical and Anesthesia Risks:}
\begin{itemize}
  \item \textbf{Anesthesia Complications:} Adverse reactions to anesthetic medications, respiratory problems (pneumonia, atelectasis), cardiac events (myocardial infarction, arrhythmias), stroke, or acute kidney injury.
  
  \item \textbf{Bleeding:} Intraoperative or postoperative hemorrhage, potentially requiring blood transfusion, re-operation, or leading to hematoma formation.
  
  \item \textbf{Infection:} Surgical site infection (SSI), ranging from superficial wound infection to deep space infection, or systemic sepsis.
  
  \item \textbf{Deep Vein Thrombosis (DVT) and Pulmonary Embolism (PE):} Blood clot formation in the deep veins of the leg, which can dislodge and travel to the lungs, potentially being life-threatening.
  
  \item \textbf{Cardiac Events:} Myocardial infarction, unstable angina, or heart failure exacerbation, particularly in patients with pre-existing coronary artery disease.
  
  \item \textbf{Stroke:} Due to embolization of plaque or thrombus, or hypoperfusion during surgery.
  
  \item \textbf{Renal Dysfunction:} Acute kidney injury, especially in patients with pre-existing renal impairment or those exposed to contrast agents.
\end{itemize}

\textbf{Procedure-Specific Risks:}
\begin{itemize}
  \item \textbf{Graft Thrombosis or Occlusion:} The most common and serious complication, where the bypass graft clots off, leading to recurrent ischemia and potentially limb loss. This often requires urgent re-intervention.
  
  \item \textbf{Graft Infection:} A devastating complication, particularly with synthetic grafts, often requiring graft removal, prolonged antibiotic therapy, and potentially leading to amputation or death.
  
  \item \textbf{Distal Embolization:} Clots or plaque fragments can dislodge during arterial clamping or flow restoration and travel downstream, blocking smaller arteries in the foot and toes, exacerbating ischemia.
  
  \item \textbf{Nerve Injury:} Damage to adjacent nerves (e.g., saphenous nerve during vein harvest, common peroneal nerve during popliteal dissection, femoral nerve during groin dissection), leading to numbness, weakness, or neuropathic pain.
  
  \item \textbf{Wound Complications:} Hematoma, seroma, lymphocele, dehiscence (wound opening), or delayed wound healing, especially in the groin or in patients with diabetes or obesity.
  
  \item \textbf{Limb Edema:} Swelling of the operated limb, which can be temporary due to lymphatic disruption or persistent due to venous hypertension.
  
  \item \textbf{Compartment Syndrome:} Rare but severe, where swelling within a muscle compartment compromises blood flow and nerve function, requiring emergency fasciotomy to prevent permanent damage.
  
  \item \textbf{Pseudoaneurysm Formation:} A false aneurysm can develop at the anastomotic site, particularly in the groin, requiring re-intervention.
  
  \item \textbf{Limb Loss:} If the bypass fails and re-intervention is not possible or unsuccessful, or if severe infection occurs, amputation may be necessary.
  
  \item \textbf{Progression of Native Arterial Disease:} The underlying atherosclerotic disease can progress in other vessels or in the native arteries proximal or distal to the bypass, leading to new ischemic symptoms or requiring further interventions.
\end{itemize}

\section*{Postoperative Care \& Recovery}
Following peripheral artery bypass surgery, patients are closely monitored in the Post-Anesthesia Care Unit (PACU) and subsequently transferred to a specialized surgical ward or intensive care unit. Immediate postoperative care focuses on ensuring graft patency and limb viability. This involves frequent assessment of vital signs, continuous monitoring of the operated limb for pulses (palpation and Doppler signals), skin temperature, color, and capillary refill. Pain management is crucial, often utilizing patient-controlled analgesia (PCA) or epidural catheters. Fluid balance is carefully managed to maintain adequate hydration and perfusion, and deep vein thrombosis prophylaxis, typically with subcutaneous heparin, is initiated.

Over the first 24-48 hours, patients are encouraged to begin ambulating with assistance to prevent complications like deep vein thrombosis and promote recovery. Antiplatelet therapy (e.g., aspirin, clopidogrel) is usually started or continued to help maintain graft patency. The surgical wounds are monitored for signs of bleeding, infection, or swelling. The hospital stay typically ranges from 3 to 7 days, depending on the complexity of the surgery, the patient's overall health, and their recovery trajectory. During the first 1-2 weeks post-discharge, patients are advised to gradually increase activity, avoid heavy lifting, and keep incisions clean and dry. Long-term recovery involves continued cardiovascular risk factor modification, including smoking cessation, diabetes control, lipid management, and blood pressure regulation, which are vital for maintaining graft patency and overall cardiovascular health. Full recovery and return to normal activities can take several weeks to months.

During peripheral artery bypass surgery, the surgeon meticulously documents the extent and characteristics of the arterial disease encountered. Typical surgical findings include severe atherosclerosis, characterized by hard, calcified plaques within the arterial wall, often leading to significant stenosis or complete occlusion. The presence of organized thrombus within the occluded segments is common. The quality of the inflow and outflow vessels, including their diameter, wall thickness, and degree of calcification, is assessed to ensure suitable anastomotic sites. The condition of the autologous vein graft, if used, is also noted, including its diameter, wall integrity, and the presence of any varicosities or previous thrombotic events.

If any arterial tissue is excised (e.g., for endarterectomy at an anastomotic site) or if a segment of occluded artery is removed, it is sent for pathological examination. The pathology report typically confirms the diagnosis of atherosclerosis, describing the presence of atherosclerotic plaques composed of lipids, inflammatory cells, smooth muscle cells, and fibrous connective tissue, often with calcification. Intimal hyperplasia, a thickening of the innermost layer of the artery, may also be noted, which is a common cause of graft stenosis. In cases of acute limb ischemia, the pathology might reveal fresh thrombus superimposed on an atherosclerotic plaque. These findings provide definitive histological confirmation of the underlying vascular disease.

A normal and successful postoperative course following peripheral artery bypass surgery is characterized by the effective restoration of blood flow to the ischemic limb and a progressive resolution of preoperative symptoms. Immediately after surgery, the operated limb should be warm, with palpable distal pulses (e.g., dorsalis pedis, posterior tibial) and normal capillary refill, indicating successful revascularization. Ischemic rest pain should resolve rapidly, often within hours or days, and non-healing ulcers or gangrene should show signs of improved perfusion, such as granulation tissue formation and reduced demarcation, leading to eventual healing over weeks to months.

Patients can expect some incisional pain, which is managed with analgesics, and mild to moderate swelling of the operated limb, which gradually subsides. Ambulation is encouraged early, typically within 24-48 hours, to prevent complications and promote recovery. Wound healing should progress without infection or dehiscence, with sutures or staples removed around 1-2 weeks post-surgery. Over the subsequent weeks, patients should experience a significant improvement in walking distance and functional capacity, allowing for a gradual return to normal daily activities. The overall goal is to achieve a durable bypass that maintains limb viability, eliminates ischemic symptoms, and improves the patient's quality of life.

Long-term follow-up is crucial after peripheral artery bypass surgery to monitor graft patency, detect recurrent disease, and manage cardiovascular risk factors effectively.

\begin{itemize}
  \item \textbf{Postoperative Visits:} Initial visits typically occur at 1-2 weeks for wound check and suture/staple removal, followed by appointments at 1 month, 3 months, 6 months, and annually thereafter with the vascular surgeon.
  
  \item \textbf{Graft Surveillance:} Lifelong surveillance with duplex ultrasound is the cornerstone of follow-up. This non-invasive imaging technique assesses graft patency, identifies stenoses within the graft or at the anastomoses, and monitors for progression of disease in native vessels. Surveillance frequency is typically higher in the first year (e.g., every 3 months) and then annually, or more frequently if concerns arise.
  
  \item \textbf{Laboratory Tests:} Regular monitoring of lipid profiles, blood glucose (HbA1c for diabetics), and renal function is important for managing cardiovascular risk factors and detecting potential complications.
  
  \item \textbf{Cardiovascular Risk Factor Management:} Continued emphasis on aggressive smoking cessation, strict control of diabetes, hypertension, and hyperlipidemia, often involving medication adjustments (e.g., statins, antiplatelets, antihypertensives) and lifestyle counseling.
  
  \item \textbf{Physical Therapy/Rehabilitation:} For patients who experienced significant deconditioning or gait abnormalities due to ischemia, a structured rehabilitation program may be beneficial to regain strength and mobility.
  
  \item \textbf{Activity Resumption:} Gradual return to normal activities is encouraged, with specific instructions on avoiding heavy lifting or strenuous activities initially, typically for 4-6 weeks.
  
  \item \textbf{Contralateral Limb Surveillance:} The contralateral limb is also at high risk for PAD progression and requires regular assessment for new or worsening symptoms.
  
  \item \textbf{Warning Signs:} Patients are educated to immediately report any new or recurrent symptoms such as rest pain, numbness, coldness, color changes in the limb, loss of pulses, or wound issues (redness, swelling, discharge) to their physician, as these may indicate graft failure or infection.
\end{itemize}

\section*{Outcomes \& Prognosis}
Peripheral artery disease (PAD) affects over 200 million people worldwide, with its prevalence increasing significantly with age, affecting up to 20\% of individuals over 60 years old. Critical limb ischemia (CLI), the most severe manifestation of PAD and the primary indication for bypass surgery, affects approximately 1-2\% of the PAD population, representing a significant burden of disease. The incidence of peripheral artery bypass surgery has seen fluctuations over the past decades, with a trend towards stabilization or slight decrease in some regions due to the rise of endovascular interventions. However, bypass remains a critical option for complex disease and limb salvage. Men tend to have a slightly higher incidence of PAD and bypass surgery than women, although the prevalence in women is increasing, particularly in older age groups. The majority of patients undergoing bypass are older adults, typically over 60 years of age, reflecting the progressive nature of atherosclerosis. Perioperative mortality for peripheral artery bypass ranges from 1\% to 5\%, depending on patient comorbidities, the urgency of the procedure, and the complexity of the bypass, while overall morbidity can be as high as 10\% to 20\%, largely driven by cardiac events, wound complications, and graft failure.

The outcomes and prognosis following peripheral artery bypass surgery are generally favorable, particularly for limb salvage, but vary significantly based on several factors including the bypass location, graft material, patient comorbidities, and adherence to postoperative care. For femoropopliteal bypasses using autologous vein, 5-year primary patency rates typically range from 60\% to 80\%, while synthetic grafts in the same position may have slightly lower rates (40-60\%). Distal bypasses (e.g., femorotibial) generally have lower patency rates but still achieve high limb salvage rates, often exceeding 70\% to 90\% at 5 years, which is the primary goal in CLI. Functional outcomes include significant improvement in walking distance, resolution of rest pain, and healing of ischemic ulcers, leading to an enhanced quality of life. Recurrence of symptoms can occur due to graft failure (thrombosis or stenosis, often due to intimal hyperplasia) or progression of atherosclerotic disease in other native vessels. Prognostic factors for better outcomes include the use of autologous vein grafts, good quality distal runoff, absence of diabetes mellitus, non-smoking status, and good overall cardiovascular health. Despite successful limb salvage, patients undergoing peripheral artery bypass remain at high risk for major adverse cardiovascular events (e.g., myocardial infarction, stroke) due to systemic atherosclerosis, underscoring the importance of aggressive long-term risk factor management. The BASIL (Bypass versus Angioplasty in Severe Ischaemia of the Leg) trial, for instance, provided valuable insights into the comparative effectiveness of bypass surgery versus endovascular therapy for critical limb ischemia, demonstrating comparable long-term outcomes but with different initial morbidity profiles.

\section*{References}
\begin{enumerate}
  \item MedlinePlus. Peripheral artery bypass surgery. U.S. National Library of Medicine; 2024. Available from: \url{https://medlineplus.gov/ency/article/002956.htm}
  
  \item Sidawy AN, Perler BA. Rutherford's Vascular Surgery and Endovascular Therapy. 10th ed. Philadelphia, PA: Elsevier; 2023.
  
  \item Society for Vascular Surgery. Clinical Practice Guidelines for the Management of Patients with Peripheral Artery Disease. J Vasc Surg. 2023;78(1S):S1-S86.
  
  \item Papadakis MA, McPhee SJ, Rabow MW. Current Medical Diagnosis and Treatment 2024. 63rd ed. New York, NY: McGraw-Hill Education; 2024.
\end{enumerate}

\clearpage
